\documentclass[../main.tex]{subfiles}
\ifSubfilesClassLoaded{\setcounter{chapter}{8}}{}
\begin{document}

\chapter{Klasifikasi: Regresi Logistik, K-NN, dan Pohon Keputusan}

\begin{subcpmk}
  \item \textbf{Sub-CPMK-P7:} Melatih minimal dua algoritma klasifikasi; menampilkan matriks konfusi pada data uji (RPS Mg.\ 8).
\end{subcpmk}

\SesiRPS{8}{P7}{$\sim$170 menit}{CPMK-P3}

\section{Kemampuan akhir sesi}

Membangun pipeline prapemrosesan sederhana; melatih klasifikasi; membandingkan akurasi atau F1; menampilkan \texttt{confusion\_matrix}.

\section{Teori}

\begin{konsep}
Klasifikasi memprediksi kelas diskrit. Regresi logistik memodelkan log-odds linear; K-NN berbasis jarak; pohon keputusan memecah ruang fitur secara hierarkis. Untuk kelas tidak seimbang, akurasi bisa menyesatkan---pertimbangkan F1 per kelas.
\end{konsep}

\section{Contoh kode}

\begin{lstlisting}[caption=Data dan split (Iris)]
from sklearn.datasets import load_iris
from sklearn.model_selection import train_test_split
X, y = load_iris(return_X_y=True)
X_train, X_test, y_train, y_test = train_test_split(
    X, y, test_size=0.25, random_state=42, stratify=y
)
\end{lstlisting}

\begin{lstlisting}[caption=LogisticRegression]
from sklearn.linear_model import LogisticRegression
log_reg = LogisticRegression(max_iter=200)
log_reg.fit(X_train, y_train)
print("logreg score", log_reg.score(X_test, y_test))
\end{lstlisting}

\begin{lstlisting}[caption=KNeighborsClassifier dan DecisionTree]
from sklearn.neighbors import KNeighborsClassifier
from sklearn.tree import DecisionTreeClassifier

knn = KNeighborsClassifier(n_neighbors=5)
knn.fit(X_train, y_train)
tree = DecisionTreeClassifier(max_depth=3, random_state=42)
tree.fit(X_train, y_train)
print("knn", knn.score(X_test, y_test))
print("tree", tree.score(X_test, y_test))
\end{lstlisting}

\begin{lstlisting}[caption=Confusion matrix]
from sklearn.metrics import confusion_matrix, ConfusionMatrixDisplay
import matplotlib.pyplot as plt
pred = log_reg.predict(X_test)
cm = confusion_matrix(y_test, pred)
disp = ConfusionMatrixDisplay(cm)
disp.plot()
plt.title("Confusion matrix (LogReg)")
plt.show()
\end{lstlisting}

\section{Referensi}

\begin{itemize}[leftmargin=*]
  \item sklearn supervised learning. \url{https://scikit-learn.org/stable/supervised_learning.html}
  \item Notebook klasifikasi SciPy 2018 (Mueller). \url{https://nbviewer.org/github/amueller/scipy-2018-sklearn/blob/master/notebooks/05.Supervised_Learning-Classification.ipynb}
\end{itemize}

\begin{asesmen}
\textbf{Tugas modul:} minimal dua model + confusion matrix + narasi model terbaik.
\end{asesmen}

\begin{rangkuman}
Minggu 8 membandingkan beberapa classifier dasar pada dataset yang sama.

\textbf{Kata kunci:} LogisticRegression, KNN, DecisionTree, confusion matrix
\end{rangkuman}

\end{document}
